% Structure of an RPC system and the steps of one call.
% Badges: -1 register, 0 bind, 1 call, 2 marshal, 3 send, 4 receive,
% 5 unmarshal, 6 actual call, 7 result.  Dashed arrows: return path.
% Usage: % Structure of an RPC system and the steps of one call.
% Badges: -1 register, 0 bind, 1 call, 2 marshal, 3 send, 4 receive,
% 5 unmarshal, 6 actual call, 7 result.  Dashed arrows: return path.
% Usage: % Structure of an RPC system and the steps of one call.
% Badges: -1 register, 0 bind, 1 call, 2 marshal, 3 send, 4 receive,
% 5 unmarshal, 6 actual call, 7 result.  Dashed arrows: return path.
% Usage: % Structure of an RPC system and the steps of one call.
% Badges: -1 register, 0 bind, 1 call, 2 marshal, 3 send, 4 receive,
% 5 unmarshal, 6 actual call, 7 result.  Dashed arrows: return path.
% Usage: \input{figures/l13-steps}   (width ~116 mm)
\begin{tikzpicture}[x=1mm, y=1mm, font=\footnotesize\sffamily,
  >={Stealth[length=1.6mm]},
  machine/.style={draw=ln-rule, rounded corners=2pt, fill=ln-box},
  comp/.style={draw, fill=white, align=center, minimum width=34mm,
               minimum height=8mm, inner sep=1pt},
  stubc/.style={comp, fill=ln-accent!15},
  rt/.style={comp, fill=ln-accent!30},
  reg/.style={draw, rounded corners=2pt, fill=white, align=center,
              minimum width=26mm, minimum height=7mm, inner sep=1pt},
  stepn/.style={circle, draw, fill=white, inner sep=0pt,
                minimum size=4.2mm, font=\scriptsize\sffamily\bfseries},
  hdr/.style={font=\scriptsize\sffamily\bfseries, text=ln-muted},
  lab/.style={font=\scriptsize\sffamily, text=ln-muted},
  ret/.style={->, densely dashed}]
  % machines
  \draw[machine] (2,0) rectangle (46,48);
  \draw[machine] (74,0) rectangle (118,48);
  \node[hdr] at (24,2.6) {\iflangja{クライアントマシン}{client machine}};
  \node[hdr] at (96,2.6) {\iflangja{サーバマシン}{server machine}};
  % client side
  \node[comp]  (caller) at (24,38)
    {\iflangja{呼び出し元}{caller}\\[-0.3ex]\texttt{s = sum(3, v);}};
  \node[stubc] (cstub)  at (24,24.5) {\iflangja{クライアントスタブ}{client stub}};
  \node[rt]    (crt)    at (24,11) {\iflangja{RPC ランタイム}{RPC runtime}};
  % server side
  \node[comp]  (proc)   at (96,38)
    {\iflangja{手続きの実装}{implementation}\\[-0.3ex]\texttt{sum(count, values)}};
  \node[stubc] (sstub)  at (96,24.5) {\iflangja{サーバスタブ}{server stub}};
  \node[rt]    (srt)    at (96,11) {\iflangja{RPC ランタイム}{RPC runtime}};
  % registry
  \node[reg] (reg) at (60,56) {\iflangja{レジストリ}{registry}};
  % -1 register, 0 bind
  \draw[->] ([xshift=-6mm]proc.north) |- (reg.east);
  \draw[<->] ([xshift=6mm]caller.north) |- (reg.west);
  \node[stepn] at (90,51.5) {$-1$};
  \node[stepn] at (30,51.5) {0};
  % client: call path down, return path up
  \draw[->]  ([xshift=-8mm]caller.south) -- ([xshift=-8mm]cstub.north);
  \draw[->]  ([xshift=-8mm]cstub.south) -- ([xshift=-8mm]crt.north);
  \draw[ret] ([xshift=8mm]crt.north) -- ([xshift=8mm]cstub.south);
  \draw[ret] ([xshift=8mm]cstub.north) -- ([xshift=8mm]caller.south);
  \node[stepn] at (11.5,31.3) {1};
  \node[stepn] at (7,24.5) {2};
  % network
  \draw[->]  ([yshift=2mm]crt.east) -- ([yshift=2mm]srt.west);
  \draw[ret] ([yshift=-2mm]srt.west) -- ([yshift=-2mm]crt.east);
  \node[stepn] at (51,17) {3};
  \node[stepn] at (69,17) {4};
  \node[lab] at (60,5.5) {\iflangja{ネットワーク}{network}};
  % server: request path up, return path down
  \draw[->]  ([xshift=-8mm]srt.north) -- ([xshift=-8mm]sstub.south);
  \draw[->]  ([xshift=-8mm]sstub.north) -- ([xshift=-8mm]proc.south);
  \draw[ret] ([xshift=8mm]proc.south) -- ([xshift=8mm]sstub.north);
  \draw[ret] ([xshift=8mm]sstub.south) -- ([xshift=8mm]srt.north);
  \node[stepn] at (113,24.5) {5};
  \node[stepn] at (83.5,31.3) {6};
  \node[stepn] at (113,38) {7};
\end{tikzpicture}
   (width ~116 mm)
\begin{tikzpicture}[x=1mm, y=1mm, font=\footnotesize\sffamily,
  >={Stealth[length=1.6mm]},
  machine/.style={draw=ln-rule, rounded corners=2pt, fill=ln-box},
  comp/.style={draw, fill=white, align=center, minimum width=34mm,
               minimum height=8mm, inner sep=1pt},
  stubc/.style={comp, fill=ln-accent!15},
  rt/.style={comp, fill=ln-accent!30},
  reg/.style={draw, rounded corners=2pt, fill=white, align=center,
              minimum width=26mm, minimum height=7mm, inner sep=1pt},
  stepn/.style={circle, draw, fill=white, inner sep=0pt,
                minimum size=4.2mm, font=\scriptsize\sffamily\bfseries},
  hdr/.style={font=\scriptsize\sffamily\bfseries, text=ln-muted},
  lab/.style={font=\scriptsize\sffamily, text=ln-muted},
  ret/.style={->, densely dashed}]
  % machines
  \draw[machine] (2,0) rectangle (46,48);
  \draw[machine] (74,0) rectangle (118,48);
  \node[hdr] at (24,2.6) {\iflangja{クライアントマシン}{client machine}};
  \node[hdr] at (96,2.6) {\iflangja{サーバマシン}{server machine}};
  % client side
  \node[comp]  (caller) at (24,38)
    {\iflangja{呼び出し元}{caller}\\[-0.3ex]\texttt{s = sum(3, v);}};
  \node[stubc] (cstub)  at (24,24.5) {\iflangja{クライアントスタブ}{client stub}};
  \node[rt]    (crt)    at (24,11) {\iflangja{RPC ランタイム}{RPC runtime}};
  % server side
  \node[comp]  (proc)   at (96,38)
    {\iflangja{手続きの実装}{implementation}\\[-0.3ex]\texttt{sum(count, values)}};
  \node[stubc] (sstub)  at (96,24.5) {\iflangja{サーバスタブ}{server stub}};
  \node[rt]    (srt)    at (96,11) {\iflangja{RPC ランタイム}{RPC runtime}};
  % registry
  \node[reg] (reg) at (60,56) {\iflangja{レジストリ}{registry}};
  % -1 register, 0 bind
  \draw[->] ([xshift=-6mm]proc.north) |- (reg.east);
  \draw[<->] ([xshift=6mm]caller.north) |- (reg.west);
  \node[stepn] at (90,51.5) {$-1$};
  \node[stepn] at (30,51.5) {0};
  % client: call path down, return path up
  \draw[->]  ([xshift=-8mm]caller.south) -- ([xshift=-8mm]cstub.north);
  \draw[->]  ([xshift=-8mm]cstub.south) -- ([xshift=-8mm]crt.north);
  \draw[ret] ([xshift=8mm]crt.north) -- ([xshift=8mm]cstub.south);
  \draw[ret] ([xshift=8mm]cstub.north) -- ([xshift=8mm]caller.south);
  \node[stepn] at (11.5,31.3) {1};
  \node[stepn] at (7,24.5) {2};
  % network
  \draw[->]  ([yshift=2mm]crt.east) -- ([yshift=2mm]srt.west);
  \draw[ret] ([yshift=-2mm]srt.west) -- ([yshift=-2mm]crt.east);
  \node[stepn] at (51,17) {3};
  \node[stepn] at (69,17) {4};
  \node[lab] at (60,5.5) {\iflangja{ネットワーク}{network}};
  % server: request path up, return path down
  \draw[->]  ([xshift=-8mm]srt.north) -- ([xshift=-8mm]sstub.south);
  \draw[->]  ([xshift=-8mm]sstub.north) -- ([xshift=-8mm]proc.south);
  \draw[ret] ([xshift=8mm]proc.south) -- ([xshift=8mm]sstub.north);
  \draw[ret] ([xshift=8mm]sstub.south) -- ([xshift=8mm]srt.north);
  \node[stepn] at (113,24.5) {5};
  \node[stepn] at (83.5,31.3) {6};
  \node[stepn] at (113,38) {7};
\end{tikzpicture}
   (width ~116 mm)
\begin{tikzpicture}[x=1mm, y=1mm, font=\footnotesize\sffamily,
  >={Stealth[length=1.6mm]},
  machine/.style={draw=ln-rule, rounded corners=2pt, fill=ln-box},
  comp/.style={draw, fill=white, align=center, minimum width=34mm,
               minimum height=8mm, inner sep=1pt},
  stubc/.style={comp, fill=ln-accent!15},
  rt/.style={comp, fill=ln-accent!30},
  reg/.style={draw, rounded corners=2pt, fill=white, align=center,
              minimum width=26mm, minimum height=7mm, inner sep=1pt},
  stepn/.style={circle, draw, fill=white, inner sep=0pt,
                minimum size=4.2mm, font=\scriptsize\sffamily\bfseries},
  hdr/.style={font=\scriptsize\sffamily\bfseries, text=ln-muted},
  lab/.style={font=\scriptsize\sffamily, text=ln-muted},
  ret/.style={->, densely dashed}]
  % machines
  \draw[machine] (2,0) rectangle (46,48);
  \draw[machine] (74,0) rectangle (118,48);
  \node[hdr] at (24,2.6) {\iflangja{クライアントマシン}{client machine}};
  \node[hdr] at (96,2.6) {\iflangja{サーバマシン}{server machine}};
  % client side
  \node[comp]  (caller) at (24,38)
    {\iflangja{呼び出し元}{caller}\\[-0.3ex]\texttt{s = sum(3, v);}};
  \node[stubc] (cstub)  at (24,24.5) {\iflangja{クライアントスタブ}{client stub}};
  \node[rt]    (crt)    at (24,11) {\iflangja{RPC ランタイム}{RPC runtime}};
  % server side
  \node[comp]  (proc)   at (96,38)
    {\iflangja{手続きの実装}{implementation}\\[-0.3ex]\texttt{sum(count, values)}};
  \node[stubc] (sstub)  at (96,24.5) {\iflangja{サーバスタブ}{server stub}};
  \node[rt]    (srt)    at (96,11) {\iflangja{RPC ランタイム}{RPC runtime}};
  % registry
  \node[reg] (reg) at (60,56) {\iflangja{レジストリ}{registry}};
  % -1 register, 0 bind
  \draw[->] ([xshift=-6mm]proc.north) |- (reg.east);
  \draw[<->] ([xshift=6mm]caller.north) |- (reg.west);
  \node[stepn] at (90,51.5) {$-1$};
  \node[stepn] at (30,51.5) {0};
  % client: call path down, return path up
  \draw[->]  ([xshift=-8mm]caller.south) -- ([xshift=-8mm]cstub.north);
  \draw[->]  ([xshift=-8mm]cstub.south) -- ([xshift=-8mm]crt.north);
  \draw[ret] ([xshift=8mm]crt.north) -- ([xshift=8mm]cstub.south);
  \draw[ret] ([xshift=8mm]cstub.north) -- ([xshift=8mm]caller.south);
  \node[stepn] at (11.5,31.3) {1};
  \node[stepn] at (7,24.5) {2};
  % network
  \draw[->]  ([yshift=2mm]crt.east) -- ([yshift=2mm]srt.west);
  \draw[ret] ([yshift=-2mm]srt.west) -- ([yshift=-2mm]crt.east);
  \node[stepn] at (51,17) {3};
  \node[stepn] at (69,17) {4};
  \node[lab] at (60,5.5) {\iflangja{ネットワーク}{network}};
  % server: request path up, return path down
  \draw[->]  ([xshift=-8mm]srt.north) -- ([xshift=-8mm]sstub.south);
  \draw[->]  ([xshift=-8mm]sstub.north) -- ([xshift=-8mm]proc.south);
  \draw[ret] ([xshift=8mm]proc.south) -- ([xshift=8mm]sstub.north);
  \draw[ret] ([xshift=8mm]sstub.south) -- ([xshift=8mm]srt.north);
  \node[stepn] at (113,24.5) {5};
  \node[stepn] at (83.5,31.3) {6};
  \node[stepn] at (113,38) {7};
\end{tikzpicture}
   (width ~116 mm)
\begin{tikzpicture}[x=1mm, y=1mm, font=\footnotesize\sffamily,
  >={Stealth[length=1.6mm]},
  machine/.style={draw=ln-rule, rounded corners=2pt, fill=ln-box},
  comp/.style={draw, fill=white, align=center, minimum width=34mm,
               minimum height=8mm, inner sep=1pt},
  stubc/.style={comp, fill=ln-accent!15},
  rt/.style={comp, fill=ln-accent!30},
  reg/.style={draw, rounded corners=2pt, fill=white, align=center,
              minimum width=26mm, minimum height=7mm, inner sep=1pt},
  stepn/.style={circle, draw, fill=white, inner sep=0pt,
                minimum size=4.2mm, font=\scriptsize\sffamily\bfseries},
  hdr/.style={font=\scriptsize\sffamily\bfseries, text=ln-muted},
  lab/.style={font=\scriptsize\sffamily, text=ln-muted},
  ret/.style={->, densely dashed}]
  % machines
  \draw[machine] (2,0) rectangle (46,48);
  \draw[machine] (74,0) rectangle (118,48);
  \node[hdr] at (24,2.6) {\iflangja{クライアントマシン}{client machine}};
  \node[hdr] at (96,2.6) {\iflangja{サーバマシン}{server machine}};
  % client side
  \node[comp]  (caller) at (24,38)
    {\iflangja{呼び出し元}{caller}\\[-0.3ex]\texttt{s = sum(3, v);}};
  \node[stubc] (cstub)  at (24,24.5) {\iflangja{クライアントスタブ}{client stub}};
  \node[rt]    (crt)    at (24,11) {\iflangja{RPC ランタイム}{RPC runtime}};
  % server side
  \node[comp]  (proc)   at (96,38)
    {\iflangja{手続きの実装}{implementation}\\[-0.3ex]\texttt{sum(count, values)}};
  \node[stubc] (sstub)  at (96,24.5) {\iflangja{サーバスタブ}{server stub}};
  \node[rt]    (srt)    at (96,11) {\iflangja{RPC ランタイム}{RPC runtime}};
  % registry
  \node[reg] (reg) at (60,56) {\iflangja{レジストリ}{registry}};
  % -1 register, 0 bind
  \draw[->] ([xshift=-6mm]proc.north) |- (reg.east);
  \draw[<->] ([xshift=6mm]caller.north) |- (reg.west);
  \node[stepn] at (90,51.5) {$-1$};
  \node[stepn] at (30,51.5) {0};
  % client: call path down, return path up
  \draw[->]  ([xshift=-8mm]caller.south) -- ([xshift=-8mm]cstub.north);
  \draw[->]  ([xshift=-8mm]cstub.south) -- ([xshift=-8mm]crt.north);
  \draw[ret] ([xshift=8mm]crt.north) -- ([xshift=8mm]cstub.south);
  \draw[ret] ([xshift=8mm]cstub.north) -- ([xshift=8mm]caller.south);
  \node[stepn] at (11.5,31.3) {1};
  \node[stepn] at (7,24.5) {2};
  % network
  \draw[->]  ([yshift=2mm]crt.east) -- ([yshift=2mm]srt.west);
  \draw[ret] ([yshift=-2mm]srt.west) -- ([yshift=-2mm]crt.east);
  \node[stepn] at (51,17) {3};
  \node[stepn] at (69,17) {4};
  \node[lab] at (60,5.5) {\iflangja{ネットワーク}{network}};
  % server: request path up, return path down
  \draw[->]  ([xshift=-8mm]srt.north) -- ([xshift=-8mm]sstub.south);
  \draw[->]  ([xshift=-8mm]sstub.north) -- ([xshift=-8mm]proc.south);
  \draw[ret] ([xshift=8mm]proc.south) -- ([xshift=8mm]sstub.north);
  \draw[ret] ([xshift=8mm]sstub.south) -- ([xshift=8mm]srt.north);
  \node[stepn] at (113,24.5) {5};
  \node[stepn] at (83.5,31.3) {6};
  \node[stepn] at (113,38) {7};
\end{tikzpicture}
